\section{Loss Function}\label{sec:method:loss}

The training objective is a sum of two \emph{supervised contrastive} \parencite{khosla2020supcon} terms with the same functional form, both aggregated over the whole mini-batch.
The two terms differ only in which label drives the positive set: an emitter identifier on the emitter branch, and the global mode label of \cref{sec:method:formulation} on the mode branch.
Routing the mode branch through a contrastive loss rather than a closed-catalogue softmax keeps the inference pipeline of \cref{eq:method:cluster_decoding} open to modulations absent from training; the post-hoc clustering operator contributes no gradients and is not part of the objective. And then embeddings for each PDW, or the encodings (before final projection), could probably be use as input to a know emitter library with softmax heads. This could also potentially enable transfer learning.

\subsection*{Supervised contrastive kernel}

All projection-head outputs of \cref{eq:method:proj_heads} are $\ell_2$-normalised, so that the cosine similarity reduces to the inner product
\begin{equation}\label{eq:method:sim}
  s_{ij} \;\coloneqq\; \vect{z}_i^{\top}\vect{z}_j \;\in\; [-1,1],
\end{equation}
and a fixed temperature $\tau > 0$ rescales it inside the loss; smaller $\tau$ concentrates gradient on the hardest negatives, larger $\tau$ stabilises early optimisation \parencite{chen2020simclr}.
Both branches use $\tau = 0.2$, the value commonly chosen for normalised contrastive features \parencite{chen2020simclr,khosla2020supcon}.

For an anchor set $I$ with embeddings $\{\vect{z}_i\}_{i \in I}$, write the candidate set $A(i) = I \setminus \{i\}$ and the positive subset $P(i) \subseteq A(i)$ defined by label coincidence.
The supervised contrastive contribution \parencite{khosla2020supcon} is
\begin{equation}\label{eq:method:supcon}
  \Loss_{\mathrm{SC}}\!\bigl(I, \{P(i)\}\bigr)
  \;=\;
  -\sum_{i \in I^{+}} \frac{1}{\lvert P(i) \rvert}
  \sum_{p \in P(i)}
  \log \frac{\exp(s_{ip}/\tau)}{\sum_{a \in A(i)} \exp(s_{ia}/\tau)},
\end{equation}
where $I^{+} = \{i \in I : P(i) \neq \emptyset\}$ and $\Loss_{\mathrm{SC}} = 0$ when $I^{+} = \emptyset$.
Anchors outside $I^{+}$ contribute only as negatives.
The denominator runs over every candidate in $A(i)$, so unrelated samples exert repulsive pressure and the encoder cannot collapse to a constant map \parencite{chen2020simclr}.

\subsection*{Branch-specific aggregation}

The two branches feed \cref{eq:method:supcon} with the same batch-wide index set
\begin{equation}\label{eq:method:batch_index}
  I \;=\; \{1,\ldots,B\}\!\times\!\{1,\ldots,L\},
\end{equation}
pooled over the $B$ windows of a mini-batch, and differ only in the labels they consume.
Let $\vect{z}^{\mathrm{em}}_{b,n}$ and $\vect{z}^{\mathrm{md}}_{b,n}$ denote the projection outputs at position $n \in \{1,\ldots,L\}$ in window $b \in \{1,\ldots,B\}$, with emitter label $e_{b,n}$ and mode label $m_{b,n}$.
Emitter identifiers are train-local in \cref{sec:method:formulation} and are taken to be globally unique within a mini-batch by pairing the train-local integer with the originating pulse-train identifier, so that $e_{b',n'} = e_{b,n}$ iff the two pulses come from the same physical emitter.
The two positive sets are
\begin{equation}\label{eq:method:positive_sets}
  P^{\mathrm{em}}_{b,n} = \{(b',n') \in I\setminus\{(b,n)\} : e_{b',n'} = e_{b,n}\},
  \qquad
  P^{\mathrm{md}}_{b,n} = \{(b',n') \in I\setminus\{(b,n)\} : m_{b',n'} = m_{b,n}\},
\end{equation}
and the two branch losses are
\begin{align}
  \Loss_{\mathrm{em}} &\;=\; \Loss_{\mathrm{SC}}\!\bigl(I,\, \{P^{\mathrm{em}}_{b,n}\}\bigr),
    \label{eq:method:loss_emitter}\\
  \Loss_{\mathrm{md}} &\;=\; \Loss_{\mathrm{SC}}\!\bigl(I,\, \{P^{\mathrm{md}}_{b,n}\}\bigr).
    \label{eq:method:loss_mode}
\end{align}
Both losses are therefore standard SupCon \parencite{khosla2020supcon} aggregated over the whole batch.
Because the loss only consumes pairwise label coincidences, neither the emitter count $K_{\mathrm{em}}$ nor the catalogue size $M$ appears in \cref{eq:method:supcon}, and modes held out at training time still inherit useful geometry from their neighbours in $\R^{q}$.
The total training objective is
\begin{equation}\label{eq:method:loss_total}
  \Loss \;=\; \Loss_{\mathrm{em}} + \lambda_{\mathrm{md}}\,\Loss_{\mathrm{md}},
  \qquad \lambda_{\mathrm{md}} > 0,
\end{equation}
with $\lambda_{\mathrm{md}}$ reported in \cref{ch:experiments}.

\subsection*{Probabilistic interpretation}

The kernel \cref{eq:method:supcon} admits a Gibbs--Boltzmann reading that clarifies the role of every factor.
For each anchor $i \in I^{+}$, define the categorical distribution over its candidates
\begin{equation}\label{eq:method:boltzmann}
  P_{\bm{\theta}}(a \mid i) \;=\; \frac{\exp(s_{ia}/\tau)}{Z_i},
  \qquad
  Z_i \;=\; \sum_{a' \in A(i)} \exp(s_{ia'}/\tau),
  \qquad a \in A(i),
\end{equation}
which is exactly the Boltzmann distribution with energy $-s_{ia}$ and temperature $\tau$ on the configuration space $A(i)$.
Substituting \cref{eq:method:boltzmann} into \cref{eq:method:supcon} yields the per-anchor cross-entropy between the empirical positive distribution (uniform on $P(i)$) and the model $P_{\bm{\theta}}(\cdot \mid i)$,
\begin{equation}\label{eq:method:supcon_logp}
  \Loss_{\mathrm{SC}}
  \;=\;
  -\sum_{i \in I^{+}} \frac{1}{\lvert P(i)\rvert}
  \sum_{p \in P(i)} \log P_{\bm{\theta}}(p \mid i),
\end{equation}
that is, minus the average log-evidence the encoder assigns to the positives, exactly as a maximum-likelihood estimator would treat $P(i)$ as a sample of observations drawn from the latent positive class.

Applying $\log(x/y) = \log x - \log y$ to \cref{eq:method:supcon} separates the attractive and repulsive contributions explicitly:
\begin{equation}\label{eq:method:supcon_split}
  \Loss_{\mathrm{SC}}
  \;=\;
  \underbrace{\sum_{i \in I^{+}} \log Z_i}_{\text{log-partition (repulsion)}}
  \;-\;
  \underbrace{\sum_{i \in I^{+}} \frac{1}{\lvert P(i)\rvert} \sum_{p \in P(i)} \frac{s_{ip}}{\tau}}_{\text{mean positive similarity (attraction)}}.
\end{equation}
The first term is monotone in similarities to every candidate, with $\partial \log Z_i / \partial s_{ia} = P_{\bm{\theta}}(a\mid i)/\tau$, so the gradient automatically concentrates on the hardest negatives---those whose Boltzmann weight under \cref{eq:method:boltzmann} is largest despite carrying a different label.
The second term rewards alignment of each anchor with its positives, with each positive contributing $1/(\tau\lvert P(i)\rvert)$ to the gradient.

The prefactor $1/\lvert P(i)\rvert$ turns the inner sum of log-probabilities into an arithmetic mean, equivalently into the log of a geometric mean,
\begin{equation}\label{eq:method:supcon_geomean}
  -\frac{1}{\lvert P(i)\rvert} \sum_{p \in P(i)} \log P_{\bm{\theta}}(p \mid i)
  \;=\;
  -\log\!\left(\,\prod_{p \in P(i)} P_{\bm{\theta}}(p \mid i)\right)^{\!1/\lvert P(i)\rvert}.
\end{equation}
Without the prefactor, the inner sum equals the log-likelihood of $\lvert P(i)\rvert$ conditionally independent observations of the latent positive class, and the loss would scale extensively with the number of available positives.
Dividing by $\lvert P(i)\rvert$ replaces this product by its geometric mean, so an anchor's contribution becomes intensive in the size of its positive set; this is what allows the same kernel to remain well-behaved across the two branches, where $\lvert P(i)\rvert$ ranges from a handful (emitter contrasts inside one window) to many hundreds (mode contrasts pooled over the batch).
